\documentclass[twocolumn]{article}
\usepackage[utf8]{inputenc}
\usepackage[T1]{fontenc}
\usepackage{amsmath, amssymb, amsfonts}
\usepackage{graphicx}
\usepackage{booktabs}
\usepackage{hyperref}

\title{Entropy-Augmented Control in Regime-Dependent Dynamical Systems}

\author{Samir Baladi\\
  \textit{Ronin Institute / Rite of Renaissance}\\
  \texttt{gitdeeper@gmail.com}
}

\date{April 2026}

\begin{document}

\maketitle

\begin{abstract}
This paper investigates the regime-dependent behavior of entropy-based control in dynamical systems. We compare PID, entropy-based nonlinear control, and a hybrid regime-switching architecture under two scenarios: an artificially unstable system and a corrected intrinsically stable system. Results demonstrate that no single controller dominates across all regimes. PID achieves optimal convergence in stable linear systems ($\Psi_{\text{final}} = 0.017$), while entropy-based control introduces mild overshoot ($\Psi_{\text{final}} = -0.239$). A hybrid switching strategy provides robust performance across regimes. These findings highlight the importance of regime-aware control design.
\end{abstract}

\section{Introduction}

The EntropyLab research program comprises ten projects. ENTRO-CORE (E-LAB-03) proposes a hybrid regime-switching control architecture combining PID with entropy-based nonlinear feedback.

\section{System Dynamics}

The entropy state $\Psi(t)$ evolves according to:

\begin{equation}
\ddot{\Psi} = -\gamma\dot{\Psi} - k\Psi + u + \xi(t)
\label{eq:dynamics}
\end{equation}

with damping $\gamma = 0.3$, restoring force $k = 0.5$, control input $u(t)$, and process noise $\xi(t) \sim \mathcal{N}(0, 0.05^2)$.

\section{Control Laws}

\subsection{PID Controller}
\begin{equation}
u_{\text{PID}}(t) = K_p e(t) + K_i \int e(\tau)d\tau + K_d \dot{e}(t)
\end{equation}
with $e(t) = -\Psi(t)$, $K_p=0.8$, $K_i=0.2$, $K_d=0.3$.

\subsection{Entropy-Based Controller}
\begin{equation}
u_{\text{ENTRO}}(t) = w_1 \sigma(\Psi_{\text{norm}} - \theta) + w_2 \tanh(\dot\Psi) + w_3 \tanh(\ddot\Psi)
\end{equation}
with $w_1=0.5$, $w_2=0.3$, $w_3=0.2$, $\theta=1.4$.

\subsection{Hybrid Regime-Switching Controller}
\begin{equation}
u(t) = \begin{cases}
u_{\text{PID}}(t), & \Psi(t) < \Psi_{\text{th}} \\
u_{\text{ENTRO}}(t), & \Psi(t) \ge \Psi_{\text{th}}
\end{cases}
\end{equation}
with switching threshold $\Psi_{\text{th}} = 1.7$.

\section{Experimental Results}

Under near-critical initial conditions ($\Psi(0)=1.8$, $\dot\Psi(0)=0.3$):

\begin{table}[h]
\centering
\caption{Terminal entropy state after 20 seconds}
\begin{tabular}{lcc}
\toprule
\textbf{Controller} & \textbf{Final $\Psi$} & \textbf{Observation} \\
\midrule
Uncontrolled & 0.053 & Naturally stable \\
PID & 0.017 & Optimal convergence \\
ENTRO-CORE & -0.239 & Mild overshoot \\
Hybrid & -0.012 & Robust performance \\
\bottomrule
\end{tabular}
\label{tab:results}
\end{table}

\section{Discussion}

Key findings:

\begin{enumerate}
\item The system is intrinsically stable due to damping ($\gamma=0.3$).
\item PID achieves the closest convergence to equilibrium ($\Psi_{\text{final}} = 0.017$).
\item ENTRO-CORE introduces mild overshoot ($\Psi_{\text{final}} = -0.239$), indicating overcorrection.
\item The hybrid controller provides robust performance without significant advantage over PID in this regime.
\end{enumerate}

\section{Conclusion}

Entropy-based control is not universally superior to PID. Its effectiveness is regime-dependent. In stable linear systems, PID remains optimal. Future work will focus on truly unstable systems where entropy-based control may provide unique benefits.

\section*{Acknowledgments}

Independent research through Ronin Institute / Rite of Renaissance.

\begin{thebibliography}{9}
\bibitem{entropia} S. Baladi, ``ENTROPIA,'' E-LAB-01, Zenodo, 2026. DOI: 10.5281/zenodo.19416737
\bibitem{entroai} S. Baladi, ``ENTRO-AI,'' E-LAB-02, Zenodo, 2026. DOI: 10.5281/zenodo.19284086
\end{thebibliography}

\end{document}
